\begin{textsample}{NEG Dim 4 – 2009\_c4 – Score -6.00 – 2009\_c4/enade\_2009\_ri\_4686cb254d.txt}  \label{ex:f4_neg_027}
Capítulo 4 Percepção sobre a Prova As percepções dos estudantes sobre a prova aplicada no ENADE/2009 na área de RELAÇÕES INTERNACIONAIS foram mensuradas por meio de nove itens que avaliaram desde o grau de dificuldade da prova até o tempo gasto para concluí -la.
As questões foram analisadas separando -se concluintes e \textbf{ingressantes}, e as percepções sobre a prova foram relacionadas com o desempenho dos estudantes e com a Região Geográfica de localização do curso6.
O desempenho dos estudantes dentro de cada Área/Subárea foi classificado em quartos ( 1/4 ) para os concluintes e para os \textbf{ingressantes} em separado ( ver a figura 1 para o esquema de representação gráfica destas grandezas ).
As notas finais foram colocadas em ordem ascendente e no primeiro quarto foram colocadas as 25 \% notas mais baixas, limitadas superiormente pelo percentil 25 que vem a ser o primeiro quartil.
No segundo quarto, as 25 \% seguintes, limitadas superiormente pelo percentil 50 que vem a ser a mediana.
No terceiro quarto, as 25 \% notas após as encontradas no segundo quarto, ou seja, as primeiras 25 \% acima da mediana, limitadas superiormente pelo percentil 75 que vem a ser o terceiro quartil.
No último quarto foram \textbf{alocadas} as 25 \% maiores notas.
Potencialmente, poderia existir uma relação entre a nota do estudante e a sua percepção sobre a prova.
Calculando as médias das respostas segundo os quartos de notas, seria possível reconhecer a existência ( ou não ) de algum \textbf{gradiente} nas respostas, por exemplo, uma melhor nota correspondendo a uma percepção menos negativa.
Este \textbf{gradiente} seria traduzido, por exemplo, no gráfico 4.1a, pela diminuição das áreas laranja e vermelha ( Difícil e Muito difícil ) com o aumento do desempenho ( indo do primeiro ao último quarto ). 6 Curso, neste relatório, refere -se à unidade de análise para o Conceito ENADE e é caracterizado pela combinação de IES, Área/Subárea ( quando pertinente ) e município de habilitação.
A medida tradicional de “ grau de discriminação ” comparou os resultados dos estudantes do primeiro quarto ( 25 \% piores estudantes ) aos do último quarto ( 25 \% melhores ).
Tal comparação, no gráfico disponibilizado, foi idêntica à comparação da primeira barra ( primeiro quarto ) com a última barra ( último quarto ) de cada grupo \textbf{ingressante} ou concluinte nos gráficos “ segundo quartos de nota final ” ( gráficos com números ímpares neste capítulo ).
Por outro lado, a disponibilização das informações segundo os quartos de nota permitiu verificar o comportamento de dois outros grupos intermediários que na análise somente dos extremos tem que ser suposto \textbf{monotônico}.
A seguir, serão apresentados os principais resultados relativos aos nove itens avaliados. 4.1 Grau de dificuldade da prova 4.1.1 Formação Geral Os Gráficos 4.1a, 4.1b e 4.1c apresentam a distribuição das respostas à questão “ Qual o grau de dificuldade desta prova na parte de Formação Geral? ”, segundo quartos de nota final e grupo de estudantes concluintes e \textbf{ingressantes} que participaram do ENADE/2009 na área de RELAÇÕES INTERNACIONAIS.
No total de concluintes, 10,7 \% declararam que a prova na parte da Formação Geral foi Difícil/Muito difícil, contra um valor próximo para os \textbf{ingressantes}, 15,3 \%.
Tanto entre os concluintes, como entre os \textbf{ingressantes}, a proporção de indivíduos que classificaram esta parte como Difícil/Muito difícil diminui com o desempenho do estudante.
Por outro lado, a proporção que declarou a parte Fácil/Muito fácil foi maior entre os concluintes ( 38,4 \% ) do que entre os \textbf{ingressantes} ( 34,0 \% ).
Observou -se um crescimento desta proporção com o aumento do desempenho entre os \textbf{ingressantes} e também entre os concluintes, que apresentam, porém, uma queda no segundo quarto.
Entre os concluintes, 50,9 \% classificaram esta parte com um nível Médio de dificuldade.
A proporção entre os \textbf{ingressantes} foi praticamente igual, 50,7 \%.
Entre os concluintes, esta proporção diminui com o desempenho, com um aumento no segundo quarto, indo de 55,1 \% no primeiro quarto de notas a 40,6 \% no último quarto.
Já entre os \textbf{ingressantes} esta proporção \textbf{decresce} com o desempenho, indo de 54,0 \% no primeiro quarto a 42,1 \% no último quarto.
Os Gráficos 4.2a e 4.2b apresentam a distribuição das respostas à questão “ Qual o grau de dificuldade desta prova na parte de Formação Geral? ”, segundo a Região onde o curso está situado para os estudantes concluintes e \textbf{ingressantes} que participaram do ENADE/2009 na área de RELAÇÕES INTERNACIONAIS.
Entre os concluintes, a maior proporção de estudantes que considerou a parte de Formação Geral da prova como de dificuldade média entre todas as Regiões brasileiras, foram os da Região Norte com 65,5 \%, já entre os \textbf{ingressantes}, foram os estudantes da Região Centro-Oeste com 52,7 \%.
Esta resposta entre os concluintes experimentou uma variação de 17 pontos percentuais com a menor proporção na Região Centro-Oeste ( 48,3 \% ).
Entre os \textbf{ingressantes} a variação para esta resposta é muito pequena, de quase 3 pontos percentuais, com a menor proporção na Região Nordeste ( 50,0 \% ). 4.1.2 Componente Específico Os Gráficos 4.3a, 4.3b e 4.3c apresentam a distribuição das respostas à questão “ Qual o grau de dificuldade desta prova na parte do Componente Específico? ”, segundo quartos de nota final e grupo de estudantes concluintes e \textbf{ingressantes} que participaram do ENADE/2009 na área de RELAÇÕES INTERNACIONAIS.
Em geral, para as desagregações consideradas, os estudantes classificaram o Componente Específico como mais difícil que a Formação Geral.
No grupo de concluintes, 29,6 \% declararam que a prova na parte do Componente Específico foi Difícil/Muito difícil, sendo um valor muito maior para os \textbf{ingressantes}, 44,7 \%.
Entre os concluintes, a proporção de indivíduos que classificaram esta parte como Muito difícil diminui com o desempenho do estudante.
A proporção que declara Difícil/Muito difícil entre os concluintes também diminui com o desempenho do estudante, com um aumento no segundo quarto, declinando 6 pontos percentuais entre o grupo de pior para o de melhor desempenho.
Entre os \textbf{ingressantes}, esta proporção cresceu entre o primeiro e o segundo quartos para então declinar até o último quarto, observando apenas os grupos extremos, o percentual declinou quase 8 pontos entre o pior e o melhor desempenho.
Entre concluintes e \textbf{ingressantes} houve estudantes que consideraram esta parte da prova como Muito fácil.
A proporção que declarou a parte Fácil/Muito fácil foi maior entre os concluintes ( 12,8 \% ) do que entre os \textbf{ingressantes} ( 6,9 \% ).
Entre os concluintes a proporção dos que consideraram esta parte como Fácil cresce a partir do segundo quarto com o desempenho.
Entre os concluintes, 57,6 \% classificaram esta parte com um nível Médio de dificuldade.
A proporção entre os \textbf{ingressantes} foi ligeiramente menor, 48,5 \%.
Entre os concluintes, esta proporção basicamente aumentou com o desempenho, indo de 50,4 \% no primeiro quarto de notas a 58,4 \% no último quarto.
Entre os \textbf{ingressantes}, esta proporção cresce a partir do segundo quarto de notas.
Os Gráficos 4.4a e 4.4b apresentam a distribuição das respostas à questão “ Qual o grau de dificuldade desta prova na parte do Componente Específico? ”, segundo a Região onde o curso está situado para os estudantes concluintes e \textbf{ingressantes} que participaram do ENADE/2009 na área de RELAÇÕES INTERNACIONAIS.
Nenhum concluinte ou \textbf{ingressante} da Região Norte declarou que esta parte foi Muito fácil ou Muito difícil.
Entre os concluintes, a maior proporção de estudantes que considerou a parte do Componente Específico da prova como de dificuldade média entre todas as Regiões brasileiras, foi na Região Norte com 65,5 \%, já entre os \textbf{ingressantes}, foram também os estudantes da Região Norte com 52,5 \%.
A Região com as menores proporções desta resposta foram respectivamente para os concluintes e \textbf{ingressantes}, a Sul ( 52,9 \% e 42,3 \% ). 4.2 Extensão da prova em relação ao tempo total Os Gráficos 4.5a, 4.5b e 4.5c apresentam a distribuição das respostas à questão sobre a extensão da prova em relação ao tempo total, segundo quartos de nota final e grupo de estudantes concluintes e \textbf{ingressantes} que participaram do ENADE/2009 na área de RELAÇÕES INTERNACIONAIS.
Tanto para concluintes quanto para \textbf{ingressantes}, a maioria considerou a extensão da prova como adequada em relação ao tempo destinado à sua resolução, respectivamente 67,2 \% e 69,1 \%.
Entre os concluintes, esta proporção cresceu com o desempenho.
Entre os \textbf{ingressantes}, esta proporção é crescente também.
A proporção dos que consideraram a extensão da prova muito curta foi ligeiramente maior entre os concluintes ( 1,2 \% ) do que entre os \textbf{ingressantes} ( 0,8 \% ).
A proporção dos que consideraram a extensão da prova curta/muito curta vis-à-vis ao tempo total foi maior entre os concluintes ( 8,6 \% ) do que entre os \textbf{ingressantes} ( 5,9 \% ).
Entre os concluintes e \textbf{ingressantes}, a declaração de que a prova teria sido muito curta com relação ao tempo total ocorre em todos os quartos e \textbf{decresce} com o desempenho entre os concluintes.
A maior proporção de concluintes e \textbf{ingressantes} declarando que a prova foi curta/muito curta foi entre os estudantes de melhor desempenho, respectivamente 11,1 \% e 10,9 \%.
A proporção dos que consideraram a extensão da prova muito longa foi maior entre os concluintes ( 7,8 \% ) do que entre os \textbf{ingressantes} ( 6,8 \% ).
A proporção dos que consideraram a extensão da prova longa foi menor entre os concluintes ( 16,4 \% ) do que entre os \textbf{ingressantes} ( 18,3 \% ).
Os Gráficos 4.6a e 4.6b apresentam a distribuição das respostas à questão sobre a extensão da prova em relação ao tempo total, segundo a Região onde o curso está situado para os estudantes concluintes e \textbf{ingressantes} que participaram do ENADE/2009 na área de RELAÇÕES INTERNACIONAIS.
Em todas as Regiões brasileiras, a maioria dos estudantes, tanto concluintes como \textbf{ingressantes}, perceberam a prova como adequada. 4.3 Compreensão dos Enunciados das Questões 4.3.1 Formação Geral Os Gráficos 4.7a, 4.7b e 4.7c apresentam a distribuição das respostas à questão “ Os enunciados das questões da prova na parte de Formação Geral, estavam claros e objetivos? ”, segundo quartos de nota final e grupo de estudantes concluintes e \textbf{ingressantes} que participaram do ENADE/2009 na área de RELAÇÕES INTERNACIONAIS.
Com relação aos enunciados das questões da parte de Formação Geral, as opiniões foram positivas : a grande maioria dos concluintes ( 81,2 \% ) e \textbf{ingressantes} ( 84,1 \% ) consideraram claros e objetivos os enunciados de todas ou da maioria das questões.
Entre os concluintes, esta proporção cresce até o terceiro quarto e então declina : 71,0 \% para o grupo com o pior desempenho, 85,8 \% para o terceiro quarto e declina para 81,7 \% no de melhor.
Entre os \textbf{ingressantes}, esta proporção apresenta o mesmo comportamento : 78,8 \% para o grupo com o pior desempenho, 88,5 \% para o terceiro grupo e 86,7 \% para o de melhor.
Os níveis de resposta informando que “ não, nenhum ” dos enunciados das questões da parte de Formação Geral era claro e objetivo foram muito baixos, tanto para concluintes ( 1,1 \% ) como para \textbf{ingressantes} ( 0,8 \% ).
Em linhas gerais, pode -se dizer que a combinação de “ poucos ” e “ nenhum ” \textbf{decresceu} com o desempenho tanto para concluintes como para \textbf{ingressantes} entre o primeiro e o terceiro quartos, experimentando uma ascensão para o último quarto ; entre os \textbf{ingressantes} a ascensão é de 3,5 \% no terceiro quarto para 4,7 \% e entre os concluintes a ascensão é de 4,2 \% no terceiro quarto para 8,8 \% no último quarto de notas.
Os Gráficos 4.8a e 4.8b apresentam a distribuição das respostas à questão “ Os enunciados das questões da prova na parte de Formação Geral, estavam claros e objetivos? ”, segundo a Região onde o curso está situado para os estudantes concluintes e \textbf{ingressantes} que participaram do ENADE/2009 na área de RELAÇÕES INTERNACIONAIS.
Em termos regionais, as proporções de estudantes que consideraram claros e objetivos os enunciados de todas ou da maioria das questões variaram entre os concluintes de 74,1 \% na Região Sul a 86,2 \% na Norte.
Já entre os \textbf{ingressantes} estas proporções variaram de 84,1 \% na Região Centro-Oeste a 96,7 \% na Região Norte.
Com relação às respostas “ poucos ” e “ não, nenhum ” tanto para concluintes como para \textbf{ingressantes} as maiores proporções ocorreram na Região Sul, respectivamente 12,3 \% e 7,4 \%. 4.3.2 Componente Específico Os Gráficos 4.9a, 4.9b e 4.9c apresentam a distribuição das respostas à questão “ Os enunciados das questões da prova na parte de formação específica estavam claros e objetivos? ”, segundo quartos de nota final e grupo de estudantes concluintes e \textbf{ingressantes} que participaram do ENADE/2009 na área de RELAÇÕES INTERNACIONAIS.
Com relação aos enunciados das questões da parte do Componente Específico, as opiniões foram também positivas : a grande maioria dos concluintes ( 78,7 \% ) e \textbf{ingressantes} ( 77,0 \% ) consideraram claros e objetivos os enunciados de todas ou da maioria das questões.
A proporção de concluintes declarando que todos os enunciados estavam claros e objetivos variou de 22,4 \% no primeiro quarto a 21,4 \% no último.
A combinação de todas ou da maioria das questões foi crescente com o desempenho para os \textbf{ingressantes}, indo de 69,7 \% para os com pior desempenho a 82,2 \% para os de melhor.
Entre os concluintes esta proporção foi crescente com o desempenho até o terceiro quarto de notas, caindo de 82,0 \% no terceiro quarto para 79,8 \% no último quarto de notas.
O número de estudantes que declararam que “ não, nenhum ” dos enunciados das questões da parte do Componente Específico era claro e objetivo, foi maior do que na parte de Formação Geral : 1,3 \% entre os concluintes e 1,0 \% entre os \textbf{ingressantes}.
Os que responderam “ poucos ” são 6,0 \% entre os concluintes e 5,8 \% entre os \textbf{ingressantes}.
A proporção de estudantes que respondeu que cerca de metade dos enunciados da parte do Componente Específico era claro e objetivo \textbf{decresceu} com o desempenho, tanto para concluintes como para \textbf{ingressantes}.
Os Gráficos 4.10a e 4.10b apresentam a distribuição das respostas à questão “ Os enunciados das questões da prova na parte de formação específica estavam claros e objetivos? ”, segundo a Região onde o curso está situado para os estudantes concluintes e \textbf{ingressantes} que participaram do ENADE/2009 na área de RELAÇÕES INTERNACIONAIS.
Em termos regionais, as proporções de \textbf{ingressantes} e de concluintes que consideraram claros e objetivos os enunciados de todas ou da maioria das questões variaram mais entre os \textbf{ingressantes} do que entre os concluintes.
Entre os concluintes, os valores foram de 76,2 \% no Sul a 89,7 \% no Norte.
Com relação às respostas “ poucos ” e “ não, nenhum ” para concluintes e \textbf{ingressantes} as maiores proporções ocorreram na Região Centro-Oeste com 8,1 \% e 8,7 \%, respectivamente. 4.4 Suficiência das informações/instruções fornecidas Os Gráficos 4.11a, 4.11b e 4.11c apresentam a distribuição das respostas à questão “ As informações/instruções fornecidas para a resolução das questões foram suficientes para resolvê -las? ”, segundo quartos de nota final e grupo de estudantes concluintes e \textbf{ingressantes} que participaram do ENADE/2009 na área de RELAÇÕES INTERNACIONAIS.
Quando perguntados sobre as informações/instruções fornecidas para a resolução das questões, quase 84,8 \% dos \textbf{ingressantes} e 89,5 \% dos concluintes de todo o Brasil responderam que estas eram ou excessivas ou suficientes em todas ou na maioria das questões, demonstrando uma avaliação positiva da prova.
A proporção de estudantes com a opinião \textbf{supracitada} é crescente com o desempenho para \textbf{ingressantes} e para concluintes é crescente até o terceiro quarto com uma queda no último.
Entre os concluintes a proporção vai de 85,8 \% para os de pior desempenho a 90,0 \% para os de melhor desempenho.
Entre os \textbf{ingressantes} as proporções são, respectivamente, 79,6 \% a 90,0 \%.
O complemento desta opinião, ou seja, os estudantes que responderam “ somente em algumas ” ou “ em nenhuma delas ”, apresenta uma proporção decrescente com o desempenho para os \textbf{ingressantes} e também para os concluintes até o terceiro quarto, com um aumento no grupo de melhor desempenho.
Os Gráficos 4.12a e 4.12b apresentam a distribuição das respostas à questão “ As informações/instruções fornecidas para a resolução das questões foram suficientes para resolvê -las? ”, segundo a Região onde o curso está situado para os estudantes concluintes e \textbf{ingressantes} que participaram do ENADE/2009 na área de RELAÇÕES INTERNACIONAIS. 4.5 Dificuldade encontrada ao responder à prova Os Gráficos 4.13a, 4.13b e 4.13c apresentam a distribuição das respostas à questão “ Você se deparou com alguma dificuldade ao responder à prova? ”, segundo quartos de nota final e grupo de estudantes concluintes e \textbf{ingressantes} que participaram do ENADE/2009 na área de RELAÇÕES INTERNACIONAIS.
Perguntados sobre as dificuldades com as quais se depararam ao responder à prova, 54,0 \% dos \textbf{ingressantes} apontaram como principal dificuldade o desconhecimento do conteúdo, já os concluintes apontaram como principal dificuldade a forma diferente de abordagem, 30,9 \%.
O desconhecimento do conteúdo foi apontado somente por 17,2 \% dos concluintes.
A forma diferente de abordagem do conteúdo foi considerada como maior dificuldade por 18,5 \% dos \textbf{ingressantes}, enquanto que a falta de motivação foi a principal dificuldade para 28,3 \% dos concluintes e para 16,7 \% dos \textbf{ingressantes}.
Os Gráficos 4.14a e 4.14b apresentam a distribuição das respostas à questão “ Você se deparou com alguma dificuldade ao responder à prova? ”, segundo a Região onde o curso está situado para os estudantes concluintes e \textbf{ingressantes} que participaram do ENADE/2009 na área de RELAÇÕES INTERNACIONAIS.
Em relação às Regiões brasileiras, o percentual de \textbf{ingressantes} que apontaram o desconhecimento do conteúdo como principal dificuldade ao responder à prova variou de 35,6 \% na Região Norte, a 56,6 \% na Sul.
Entre os concluintes, tal percentual variou de 17,0 \% na Região Sudeste a 17,9 \%, na Sul.
A forma diferente de abordagem de conteúdo foi citada como maior dificuldade por 21,5 \% dos concluintes na Região Nordeste chegando a 48,3 \% na Região Norte.
Já entre os \textbf{ingressantes} os valores foram de 17,2 \% na Região Sudeste a 23,7 \% na Região Norte.
O percentual de concluintes que citaram a falta de motivação como principal dificuldade variou de 10,3 \% na Região Norte a 38,3 \% na Região Nordeste.
Entre os \textbf{ingressantes}, as estatísticas correspondentes foram de 13,2 \% na Região Sul a 20,3 \%, na Região Norte.
Os que declararam que não tiveram qualquer tipo de dificuldade variaram, entre os concluintes, de 13,8 \% na Região Centro-Oeste a 20,7 \% na Região Norte.
Entre os \textbf{ingressantes}, este percentual variou de 6,0 \% na Região Sul a 11,9 \% na Região Norte. 4.6 Percepção sobre os conteúdos das questões objetivas da prova Os Gráficos 4.15a, 4.15b e 4.15c apresentam a distribuição das respostas à questão “ Considerando apenas as questões objetivas da prova, você percebeu que :... ”, segundo quartos de nota final e grupo de estudantes concluintes e \textbf{ingressantes} que participaram do ENADE/2009 na área de RELAÇÕES INTERNACIONAIS.
As respostas possíveis para esta questão estavam relacionadas à cobertura do conteúdo das questões objetivas durante o curso.
A proporção de \textbf{ingressantes} ( 57,3 \% ) que declarou que não estudou ainda a maioria desses conteúdos foi bem maior do que entre os concluintes ( 2,8 \% ).
Esta proporção foi decrescente tanto para os concluintes como para os \textbf{ingressantes} em função do desempenho, embora entre estes últimos aconteça um pequeno aumento do primeiro quarto para o segundo.
Por outro lado, a proporção de estudantes que declarou que estudou e aprendeu todos ou a maioria desses conteúdos, entre os \textbf{ingressantes} foi de 25,5 \% e entre os concluintes 75,2 \%.
Esta proporção foi crescente em função do desempenho, tanto para os concluintes como para os \textbf{ingressantes}.
A proporção de estudantes que declarou ter estudado a maioria desses conteúdos, mas não os ter aprendido foi bem maior entre os concluintes ( 15,1 \% ) do que os \textbf{ingressantes} ( 5,7 \% ).
Os Gráficos 4.16a e 4.16b apresentam a distribuição das respostas à questão “ Considerando apenas as questões objetivas da prova, você percebeu que :... ”, segundo a Região onde o curso está situado para os estudantes concluintes e \textbf{ingressantes} que participaram do ENADE/2009 na área de RELAÇÕES INTERNACIONAIS.
Em relação às Regiões brasileiras, concluintes que afirmaram não ter estudado ainda a maioria dos conteúdos ocorreram em todas as Regiões, variando de 2,3 \% na Região Centro-Oeste a 3,7 \% na Região Nordeste.
Entre os \textbf{ingressantes}, os percentuais variaram de 55,8 \% na Região Sudeste a 62,8 \%, na Centro-Oeste.
Em todas as Regiões, a maioria dos concluintes afirmou ter estudado e aprendido todos ou muitos dos conteúdos, com percentuais que variaram de 68,6 \% na Região Sul a 76,5 \%, na Sudeste.
Para os \textbf{ingressantes}, tal percentual variou de 22,6 \%, na Região Sul, a 30,5 \%, na Norte. 4.7 Tempo gasto para concluir a prova Os Gráficos 4.17a, 4.17b e 4.17c apresentam a distribuição das respostas à questão “ Qual foi o tempo gasto para concluir a prova? ”, segundo quartos de nota final e grupo de estudantes concluintes e \textbf{ingressantes} que participaram do ENADE/2009 na área de RELAÇÕES INTERNACIONAIS.
Quando perguntados quanto ao tempo de conclusão da prova, a maioria respondeu entre uma e duas horas, tanto entre os concluintes ( 46,6 \% ) quanto entre os \textbf{ingressantes} ( 48,2 \% ).
A segunda maior incidência de respostas foi entre duas e três horas, igualmente para ambos os grupos : concluintes ( 33,4 \% ) e \textbf{ingressantes} ( 32,5 \% ).
A proporção daqueles que declararam ter usado as quatro horas, mas não ter conseguido terminar foi bem baixa : 0,9 \% entre os \textbf{ingressantes} e 1,0 \% entre os concluintes.
Entre os concluintes, esta resposta ficou concentrada entre os estudantes do primeiro quarto, aquele com pior desempenho.
Nem entre os concluintes, nem entre os \textbf{ingressantes} percebeu -se uma tendência \textbf{monotônica} do tempo necessário para concluir a prova como função do desempenho.
Tanto entre os \textbf{ingressantes} quanto entre os concluintes os grupos de mais baixo desempenho, aqueles que levam mais tempo para concluir a prova, estão no primeiro quarto de notas.
Os Gráficos 4.18a e 4.18b apresentam a distribuição das respostas à questão “ Qual foi o tempo gasto para concluir a prova? ”, segundo a Região onde o curso está situado para os estudantes concluintes e \textbf{ingressantes} que participaram do ENADE/2009 na área de RELAÇÕES INTERNACIONAIS.
A proporção de concluintes que gastaram entre uma e duas horas foi maior na Região Nordeste ( 51,8 \% ) e menor na Região Norte ( 35,7 \% ).
Para \textbf{ingressantes}, o maior percentual para este tempo ocorreu na Região Sul ( 50,6 \% ) e o menor na Região Norte ( 41,7 \% ).
Considerando -se a distribuição cumulativa, entre os concluintes, os estudantes da Região Sul foram os que levaram mais tempo para concluir a prova por oposição aos estudantes da Região Norte que levaram menos tempo.
Entre os \textbf{ingressantes}, as Regiões correspondentes foram a Nordeste com tempo mais longo e a Norte com tempo mais curto.

% matched lemmas: alocar, decrescer, gradiente, ingressante, monotônico, supracitado
\end{textsample}
